\documentclass{article}
\usepackage[utf8]{inputenc}
\usepackage[T1]{fontenc}
\usepackage{amsmath,amssymb,amsthm}
\usepackage{booktabs}
\usepackage{geometry}
\usepackage{hyperref}

\hypersetup{
  colorlinks=true,
  linkcolor=blue,
  citecolor=blue,
  urlcolor=blue,
  pdftitle={H$_3$ → E$_8$ → SU(3) Projection and Golden Ratio},
  pdfauthor={Dmitrii Vasilev}
}

\title{H$_3$ → E$_8$ → SU(3) Projection Analysis: Does \varphi$ Appear in SU(3) Casimir Operators?}
\author{Dmitrii Vasilev}
\date{2026-04-13}

\begin{document}
\maketitle

\begin{abstract}
The Coxeter chain $H_3 \xrightarrow{\text{spinors}} H_4 \xrightarrow{\text{McKay}} E_8$ provides a geometric pathway for the golden ratio $\varphi$ to enter particle physics. The $E_8$ exceptional Lie group contains $\varphi$ as a structural constant through its Coxeter number $h(E_8) = 30$ and through the subgroup chain $H_3 \subset H_4 \subset E_8$. We analyze whether this geometric connection propagates to the SU(3) color gauge group embedded in $E_8$, specifically examining whether $\varphi$ appears in the quadratic Casimir operator $C_2$ or Dynkin index $T$ of the relevant SU(3) representations.
\end{abstract}

\section{Introduction}

\subsection{The Golden Ratio in $E_8$ and $H_3$, $H_4$}

The exceptional group $E_8$ has the following golden-ratio properties~\cite{Baez2002}:

\begin{enumerate}
  \item \textbf{Coxeter number:} $h(E_8) = 30 = \varphi + \tau$ where $\tau = \varphi$ (since $\tau = (1+\sqrt{5})/2$).
  \item \textbf{Root system:} $E_8$ has 240 roots in 8-dimensional space with coordinates in $\mathbb{Q}(\sqrt{5})$.
  \item \textbf{$H_4$ subgroup:} The 120-root $H_4$ icosahedral group (4D) has all roots in $\mathbb{Z}[\varphi]$.
  \item \textbf{Geometric $\varphi$:} The angle between root vectors in $H_4$ involves $\arccos(\varphi/2) \approx 36^\circ$.
\end{enumerate}

\subsection{Coxeter Chain Connection}

\begin{equation}
  H_3 \xrightarrow{\text{double cover}} \tilde{H}_3 \xrightarrow{\text{spinors}} H_4 \xrightarrow{\text{McKay}} E_8
  \label{eq:coxeter-chain}
\end{equation}

Where:
\begin{itemize}
  \item $H_3$: Icosahedral group in 3D (15 roots)
  \item $\tilde{H}_3$: Double cover of $H_3$ in 4D
  \item $H_4$: Icosahedral group in 4D (120 roots)
  \item $E_8$: Exceptional Lie group of dimension 248
\end{itemize}

\subsection{The McKay Correspondence}

The McKay correspondence relates finite subgroups of $E_8$ to Dynkin diagrams of Lie groups. A key pathway is:

\begin{equation}
  H_4 \xrightarrow{\text{McKay}} E_8 \supset SO(3) \times SO(5) \times G_2 \supset \text{SU}(3) \times E_6
  \label{eq:mckay}
\end{equation}

Where $\text{SU}(3)$ appears as a maximal subgroup on multiple paths.

\section{$E_8 \to \text{SU}(3)$ Branching Rules}

\subsection{Maximal SU(3) Subgroup in $E_8$}

The $E_8$ root system contains $\text{SU}(3) \times E_6$ as a maximal subgroup. The branching rules are:

\begin{table}[h]
\centering
\caption{Branching of selected $E_8$ representations to SU(3) multiplets.}
\label{tab:e8-su3-branching}
\renewcommand{\arraystretch}{1.2}
\begin{tabular}{lcc}
\toprule
$E_8$ Irrep & Dimension & SU(3) Branching \\
\midrule
$\mathbf{248}$ (adjoint) & 248 & $3 + 6 + 15 + 3$ \\
$\mathbf{248}'$ (fundamental) & 248' & $3 \times 82$ \\
$248 \oplus 248'$ & 384 & Complex decomposition \\
\bottomrule
\end{tabular}
\end{table}

\subsection{Casimir Invariants of SU(3) Representations}

For SU(3) representations with Dynkin label $(p, q)$, the quadratic Casimir is:

\begin{equation}
  C_2(p, q) = \frac{p^2 + q^2 + pq + 3p + 3q}{3}
  \label{eq:su3-c2}
\end{equation}

Key values:

\begin{table}[h]
\centering
\caption{Quadratic Casimir $C_2$ for SU(3) representations.}
\label{tab:su3-c2}
\renewcommand{\arraystretch}{1.1}
\begin{tabular}{lccc}
\toprule
Name & Dynkin $(p, q)$ & Dimension & $C_2$ \\
\midrule
Fundamental (3) & $(1, 0)$ & 3 & $4/3$ \\
Adjoint (8) & $(1, 1)$ & 8 & $3$ \\
Decuplet (10) & $(3, 0)$ & 10 & $6$ \\
27-plet (27) & $(2, 2)$ & 27 & $10$ \\
\bottomrule
\end{tabular}

\section{Projection Analysis: Does \varphi$ Propagate?}

\subsection{Geometric $\varphi$ in $E_8$ Root Coordinates}

The $E_8$ root system contains explicit $\varphi$ in its coordinate representation. A representative root pair is:

\begin{equation}
  \alpha_1 = (\varphi, \varphi, -1/\varphi, \varphi, \dots)
  \label{eq:e8-phi-root}
\end{equation}

where the coordinates involve $\varphi = (1+\sqrt{5})/2 \approx 1.618$.

\subsection{Test: Projection to SU(3) Casimir}

The question is whether the quadratic Casimir operator $C_2$ for any SU(3) representation contains $\varphi$.

\textbf{Hypothesis:} If $\alpha_s = \varphi^{-3}/2$ is a group-theoretic invariant, then for some SU(3) representation $\mathbf{R}$ appearing in the $E_8 \to \text{SU}(3)$ branching, we should have:

\begin{equation}
  \frac{C_2(\mathbf{R}_\text{phys})}{T(\mathbf{R}_\text{phys})} \stackrel{?}{=} \varphi^{-3/2}
  \label{eq:hypothesis}
\end{equation}

where $\mathbf{R}_\text{phys}$ is a physical SU(3) representation (adjoint or fundamental) relevant to QCD.

\subsection{Analysis of Specific Branching Cases}

\paragraph{Case 1: Adjoint $\mathbf{8}$ from $\mathbf{248}$}

The adjoint representation $\mathbf{8}$ has $C_2 = 3$ and $T = 8$.

\begin{itemize}
  \item $\mathbf{R} = \mathbf{8}$: This is the QCD color gauge group.
  \item Ratio: $C_2/T = 3/8 = 0.375$.
  \item Comparison to $\varphi^{-3/2} \approx 0.118$: No match.
  \item Geometric $\varphi$ from Eq.~(\ref{eq:e8-phi-root}) enters through $E_8$ root coordinates but does \emph{not} affect the algebraic Casimir value.
\end{itemize}

\paragraph{Case 2: Fundamental $\mathbf{3}$ from $\mathbf{248}'$}

The fundamental representation $\mathbf{3}$ has $C_2 = 4/3$ and $T = 1$.

\begin{itemize}
  \item $\mathbf{R} = \mathbf{3}$: This is a quark flavor triplet representation.
  \item Ratio: $C_2/T = (4/3)/1 = 4/3 \approx 1.333$.
  \item Comparison to $\varphi^{-3/2}$: No match.
\end{itemize}

\paragraph{Case 3: Mixed Representations}

For branched representations involving $\mathbf{3}$ and $\mathbf{8}$:

\begin{itemize}
  \item $\mathbf{R} = \mathbf{3} \oplus \mathbf{8}$ (total dimension 11): Mixed Casimir involves tensor products.
  \item The Casimir of a tensor product is not a simple ratio of pure representations.
  \item No simple $\varphi$-dependence emerges from these mixtures.
\end{itemize}

\subsection{Geometric vs. Algebraic Transmission}

The geometric $\varphi$ in $E_8$ coordinates is a \emph{spatial} feature of the root system. However, the Casimir operator $C_2$ is an \emph{algebraic} invariant depending only on the Lie algebra structure.

\textbf{Key insight:} The golden ratio in $E_8$ geometry is \emph{coordinate-dependent}---it appears in specific root vector coordinates---while $C_2$ is \emph{coordinate-independent}---it is the same for all isomorphic representations.

% \begin{theorem}[Invariant Independence]
   The quadratic Casimir $C_2$ of an SU(3) representation is invariant under changes of basis, and therefore cannot depend on the geometric choice of $E_8$ root coordinates.
\end{theorem}
  The quadratic Casimir $C_2$ of an SU(3) representation is invariant under changes of basis, and therefore cannot depend on the geometric choice of $E_8$ root coordinates.
\end{theorem}

\subsection{Fourth-Order Invariant $C_4$}

For completeness, we examine the fourth-order Casimir:

\begin{equation}
  C_4(p, q) = \frac{2(p^4 + q^4) + 4(p^3q + pq^3) + 6(p^2q^2 + pq^2) + 12(p^3 + q^3) + 36pq + 27(p^2 + q^2)}{27}
  \label{eq:su3-c4}
\end{equation}

For fundamental $\mathbf{3}$:
\begin{equation}
  C_4(1, 0) = \frac{2(1 + 0) + 4(1 \cdot 0 + 0) + 6(1 \cdot 0^2 + 1 \cdot 0) + 12(1^3 + 0^3) + 36(1 \cdot 0) + 27(1^2 + 0^2)}{27}
         = \frac{2 + 0 + 0 + 0 + 0 + 12 + 0 + 0 + 27}{27}
         = \frac{41}{27} \approx 1.518
\end{equation}

No $\varphi$ dependence emerges. For adjoint $\mathbf{8}$, $C_4$ involves $p=1, q=1$ and evaluates to the same structure.

\subsection{Alternative: Group Volume Ratios}

The volume of a Lie group in root space provides an alternative invariant. For SU(3):

\begin{equation}
  \text{Vol}(\text{SU}(3)) = \frac{4\pi^{3/2}}{\sqrt{3}} \approx 19.076
  \label{eq:su3-volume}
\end{equation}

Comparison:
\begin{itemize}
  \item $\varphi^3 \approx 4.236$
  \item $\text{Vol}/\pi^2 \approx 19.076/\pi^2 \approx 1.934$
  \item No equality: $\varphi^3 \neq \text{Vol}/\pi^2$
\end{itemize}

\section{Higher-Dimensional Coxeter Groups}

\subsection{$H_4$ $\to$ SU(3) Embedding Question}

The $H_4$ icosahedral group contains $\text{SU}(3)$ as a subgroup? The maximal subgroups of $H_4$ are:

\begin{itemize}
  \item $A_4$: Alternating group on 4 elements (order 12)
  \item $D_4$: Dihedral group of order 24
  \item $\text{SU}(3)$: \textbf{Not} a subgroup of $H_4$
\end{itemize}

This is significant: $H_4$'s geometric $\varphi$ does \emph{not} connect to SU(3) because SU(3) is not embedded in $H_4$.

\subsection{Direct $H_3 \to \text{SU}(3)$?}

The 3D icosahedral group $H_3$ has Weyl group $A_5$ as double cover. The $A_5$ subgroup of $E_8$ (via McKay) maps to SU(3) in certain embeddings. However:

\begin{itemize}
  \item $H_3$ is not a subgroup of $E_8$
  \item The McKay path goes $H_3 \to \tilde{H}_3 \to H_4 \to E_8 \supset \text{SU}(3)$
  \item The geometric $\varphi$ appears at $H_4$ level (in $\mathbb{Z}[\varphi]$ root coordinates)
  \item This $\varphi$ is \emph{lost} when mapping to SU(3) because SU(3) has discrete roots at $60^\circ$ angles, not $36^\circ$
\end{itemize}

\section{Conclusion}

\subsection{Summary of Findings}

\begin{enumerate}
  \item \textbf{Geometric $\varphi$ in $E_8$:} The golden ratio appears in $E_8$ root coordinates, but this is a coordinate-dependent feature.
  \item \textbf{SU(3) invariance:} The quadratic Casimir $C_2$ is coordinate-independent and does not contain $\varphi$ for any SU(3) representation.
  \item \textbf{No propagation:} The $E_8 \to \text{SU}(3)$ branching does \emph{not} transmit geometric $\varphi$ to SU(3) Casimir operators.
  \item \textbf{Alternative mechanisms:} Group volume ratios and higher-order invariants $C_4$ also show no $\varphi$ connection.
  \item \textbf{McKay correspondence:} Provides a structural link between $H_4$ and $E_8$, but SU(3) appears as an independent subgroup, not through a $\varphi$-preserving channel.
\end{enumerate}

\subsection{Implication for $\alpha_s$}

\textbf{Primary finding:} The geometric pathway $H_3 \xrightarrow{H_4} \xrightarrow{E_8} \xrightarrow{\text{SU}(3)}$ does not provide a mechanism linking $\varphi$ to $\alpha_s = \varphi^{-3}/2$.

\textbf{Explanation:} The golden ratio in $E_8$ is embedded in the root system geometry, while QCD's $\alpha_s$ is determined by SU(3) algebraic invariants. These are distinct mathematical structures, and the $\varphi$ in $E_8$ geometry does not influence SU(3) Casimir values.

This analysis strengthens the conclusion of the main $\alpha_s$ paper: the coincidence $\alpha_s(m_Z) \approx \varphi^{-3}/2$ has no known theoretical foundation in SU(3) group theory or $E_8$ exceptional group geometry.

\section{Alternative: A$_5$ Discrete Symmetry Path}

The null result $\sin^2\theta_{23} - \varphi^{-4} = 0$ from Trinity suggests a different discrete symmetry connection. A$_5$ has Coxeter number $h(A_5) = 5$ (not $\tau$), but its representations involve the same 5th roots of unity that generate $\varphi$.

The A$_5$ $\to \text{SU}(3)$ branching rules are analyzed separately (see \texttt{a5\_su3\_branching.tex}).

\begin{thebibliography}{99}

\bibitem{Baez2002}
J.~C. Baez,
\textit{The Octonions},
\textit{Bull.\ Amer.\ Math.\ Soc.} \textbf{39}, 145--205 (2002).
\href{https://doi.org/10.1090/S0273-0979-01-00934-X}{DOI:~10.1090/S0273-0979-01-00934-X}

\bibitem{McKay1980}
J.~McKay,
\textit{Graphs, Singularities, and Finite Groups},
\textit{Bull.\ London\ Math.\ Soc.} \textbf{43}, 437--452 (1980).

\bibitem{GrossWilczek1973}
D.~J.~Gross and F.~Wilczek,
\textit{Ultraviolet Behavior of Non-Abelian Gauge Theories},
\textit{Phys.\ Rev.\ Lett.} \textbf{30}, 1343--1346 (1973).

\end{thebibliography}

\end{document}
